% Options for packages loaded elsewhere
\PassOptionsToPackage{unicode}{hyperref}
\PassOptionsToPackage{hyphens}{url}
%
\documentclass[
  11pt,
]{article}
\usepackage{amsmath,amssymb}
\usepackage{iftex}
\ifPDFTeX
  \usepackage[T1]{fontenc}
  \usepackage[utf8]{inputenc}
  \usepackage{textcomp} % provide euro and other symbols
\else % if luatex or xetex
  \usepackage{unicode-math} % this also loads fontspec
  \defaultfontfeatures{Scale=MatchLowercase}
  \defaultfontfeatures[\rmfamily]{Ligatures=TeX,Scale=1}
\fi
\usepackage{lmodern}
\ifPDFTeX\else
  % xetex/luatex font selection
  % Native editor uses its default monospaced font.
\fi
% Use upquote if available, for straight quotes in verbatim environments
\IfFileExists{upquote.sty}{\usepackage{upquote}}{}
\IfFileExists{microtype.sty}{% use microtype if available
  \usepackage[]{microtype}
  \UseMicrotypeSet[protrusion]{basicmath} % disable protrusion for tt fonts
}{}
\makeatletter
\@ifundefined{KOMAClassName}{% if non-KOMA class
  \IfFileExists{parskip.sty}{%
    \usepackage{parskip}
  }{% else
    \setlength{\parindent}{0pt}
    \setlength{\parskip}{6pt plus 2pt minus 1pt}}
}{% if KOMA class
  \KOMAoptions{parskip=half}}
\makeatother
\usepackage{xcolor}
\usepackage[margin=1in]{geometry}
\usepackage{color}
\usepackage{fancyvrb}
\newcommand{\VerbBar}{|}
\newcommand{\VERB}{\Verb[commandchars=\\\{\}]}
\DefineVerbatimEnvironment{Highlighting}{Verbatim}{commandchars=\\\{\}}
% Add ',fontsize=\small' for more characters per line
\newenvironment{Shaded}{}{}
\newcommand{\AlertTok}[1]{\textcolor[rgb]{1.00,0.00,0.00}{\textbf{#1}}}
\newcommand{\AnnotationTok}[1]{\textcolor[rgb]{0.38,0.63,0.69}{\textbf{\textit{#1}}}}
\newcommand{\AttributeTok}[1]{\textcolor[rgb]{0.49,0.56,0.16}{#1}}
\newcommand{\BaseNTok}[1]{\textcolor[rgb]{0.25,0.63,0.44}{#1}}
\newcommand{\BuiltInTok}[1]{\textcolor[rgb]{0.00,0.50,0.00}{#1}}
\newcommand{\CharTok}[1]{\textcolor[rgb]{0.25,0.44,0.63}{#1}}
\newcommand{\CommentTok}[1]{\textcolor[rgb]{0.38,0.63,0.69}{\textit{#1}}}
\newcommand{\CommentVarTok}[1]{\textcolor[rgb]{0.38,0.63,0.69}{\textbf{\textit{#1}}}}
\newcommand{\ConstantTok}[1]{\textcolor[rgb]{0.53,0.00,0.00}{#1}}
\newcommand{\ControlFlowTok}[1]{\textcolor[rgb]{0.00,0.44,0.13}{\textbf{#1}}}
\newcommand{\DataTypeTok}[1]{\textcolor[rgb]{0.56,0.13,0.00}{#1}}
\newcommand{\DecValTok}[1]{\textcolor[rgb]{0.25,0.63,0.44}{#1}}
\newcommand{\DocumentationTok}[1]{\textcolor[rgb]{0.73,0.13,0.13}{\textit{#1}}}
\newcommand{\ErrorTok}[1]{\textcolor[rgb]{1.00,0.00,0.00}{\textbf{#1}}}
\newcommand{\ExtensionTok}[1]{#1}
\newcommand{\FloatTok}[1]{\textcolor[rgb]{0.25,0.63,0.44}{#1}}
\newcommand{\FunctionTok}[1]{\textcolor[rgb]{0.02,0.16,0.49}{#1}}
\newcommand{\ImportTok}[1]{\textcolor[rgb]{0.00,0.50,0.00}{\textbf{#1}}}
\newcommand{\InformationTok}[1]{\textcolor[rgb]{0.38,0.63,0.69}{\textbf{\textit{#1}}}}
\newcommand{\KeywordTok}[1]{\textcolor[rgb]{0.00,0.44,0.13}{\textbf{#1}}}
\newcommand{\NormalTok}[1]{#1}
\newcommand{\OperatorTok}[1]{\textcolor[rgb]{0.40,0.40,0.40}{#1}}
\newcommand{\OtherTok}[1]{\textcolor[rgb]{0.00,0.44,0.13}{#1}}
\newcommand{\PreprocessorTok}[1]{\textcolor[rgb]{0.74,0.48,0.00}{#1}}
\newcommand{\RegionMarkerTok}[1]{#1}
\newcommand{\SpecialCharTok}[1]{\textcolor[rgb]{0.25,0.44,0.63}{#1}}
\newcommand{\SpecialStringTok}[1]{\textcolor[rgb]{0.73,0.40,0.53}{#1}}
\newcommand{\StringTok}[1]{\textcolor[rgb]{0.25,0.44,0.63}{#1}}
\newcommand{\VariableTok}[1]{\textcolor[rgb]{0.10,0.09,0.49}{#1}}
\newcommand{\VerbatimStringTok}[1]{\textcolor[rgb]{0.25,0.44,0.63}{#1}}
\newcommand{\WarningTok}[1]{\textcolor[rgb]{0.38,0.63,0.69}{\textbf{\textit{#1}}}}
\setlength{\emergencystretch}{3em} % prevent overfull lines
\providecommand{\tightlist}{%
  \setlength{\itemsep}{0pt}\setlength{\parskip}{0pt}}
\setcounter{secnumdepth}{-\maxdimen} % remove section numbering
\usepackage{fvextra}
\RecustomVerbatimEnvironment{Highlighting}{Verbatim}{commandchars=\\\{\},breaklines=true,fontsize=\footnotesize}
\ifLuaTeX
  \usepackage{selnolig}  % disable illegal ligatures
\fi
\usepackage{bookmark}
\IfFileExists{xurl.sty}{\usepackage{xurl}}{} % add URL line breaks if available
\urlstyle{same}
\hypersetup{
  pdftitle={Contact for every fixed real reference on products of rank-one balanced graphs},
  pdfauthor={Xean research manuscript},
  hidelinks,
  pdfcreator={LaTeX via pandoc}}

\title{Contact for every fixed real reference on products of rank-one
balanced graphs}
\author{Xean research manuscript}
\date{4 October 2026}

\begin{document}
\maketitle

\subsection{Exact new theorem and residual
scope}\label{exact-new-theorem-and-residual-scope}

Let \(G_i\subset N_{i,\mathbb{R}}\), \(i=1,2\), be nonempty finite
embedded balanced rational pure one-dimensional polyhedral spaces, in
their \textbf{actual lattices}, with arbitrary strictly positive
\textbf{real} edge weights. For every vertex, including bivalent
subdivisions, require the actual saturated primitive outgoing directions
\(v_e\) and actual weights \(w_e\) to satisfy \[
 \{c\geq0:\sum_e c_e v_e=0\}=\mathbb{R}_{\geq0}(w_e)_e,
 \qquad w_e>0.                                                \tag{1}
\] Retain all valences, bounded edges, cycles, ends, full lines, and
finitely many disconnected components. Give \(X=G_1\times G_2\) its
actual product lattice and product weights. Call a factor
\textbf{unbranched} when every carrier junction is bivalent, allowing a
structure consisting of a whole line.

Fix \textbf{each prescribed} continuous real-coefficient PA actual-PC
reference \(\gamma\) on a finite \textbf{rational spatial} structure,
positive at infinity. All function and operator conventions here are the
explicit real scalar extension on finite rational domains; rational
spatial domains are not inferred from real affine comparison
coefficients.

\textbf{Fixed-real-reference theorem.} For this fixed
pair \((X,\gamma)\), every bounded continuous real-PA rational-domain
obstacle has a finite continuous rational-domain actual-PAPC raw
envelope in \(\gamma+O(1)\), uniformly decreasing PAPC regularization,
and complete contact \[
 \int_X(P_\gamma(u)-\gamma-u)
                \mathrm{MA}_{\mathrm{poly},\mathbb{R}}(P_\gamma(u))=0    \tag{2}
\] \textbf{if and only if at least one factor is unbranched}. If both
branch, for \textbf{every} such prescribed real \(\gamma\) and any
selected branched junctions \(A,B\), there is a nonnegative compactly
supported bounded continuous real-PA obstacle on finite rational spatial
domains, with actual PAPC regularizable exact raw envelope \(P\),
unchanged reference at infinity, and the \textbf{full} strictly negative
integral \[
 P_\gamma(u)=P,\qquad
 \int_X(P-\gamma-u)\,\mathrm{MA}_{\mathrm{poly},\mathbb{R}}(P)
                             =-8\varepsilon^2ab<0.            \tag{3}
\] Here \(a,b>0\) are the actual real scales of the primitive positive
integer star relations, and \(\varepsilon>0\) is real. The operator is
\(P\cdot(P\cdot[X])\), with \textbf{no division by \(2!\)}.

The new step is an \textbf{exact finite convex decomposition} of a
prescribed real reference into rational positive-growth actual-PC
references on the same fixed complex, followed by a positive
\textbf{sum} of their rational-domain localizations. It removes the
rational-reference restriction in this two-graph subclass. The earlier
rational case and positive cylinder argument are supporting mathematics
reconstructed below, not new full-atom results or theorem grants.

The original P05.2 target is: characterize the pairs for which the
regularized polyhedral envelope has the exact contact identity for all
admissible obstacles. The exact full original problem and completion
criteria are preserved verbatim as strings in \texttt{task.json} and in
\texttt{original-scope.txt}, extracted from \texttt{input.json}; the
current assignment is in \texttt{assignment.txt}. These files are
mathematical scope data. The criteria require a rigorous complete proof
or an analytic counterexample satisfying every ambient hypothesis for
each claimed solution, and explicitly forbid treating a known case,
extra conditional hypothesis, abstract tropical obstruction, or narrower
theorem as full closure. A new special case can be reported with its
exact scope and the full atom left open.

This is a reviewed partial theorem. It does not
classify arbitrary balanced spaces, other higher-dimensional carriers or
products, additional affine factors, the intrinsic smoothness and
strict-positivity contact package, literal coefficients, analytic
realizations, or P06.2. No all-balanced-space regularity or literal
closure is inferred from the fixed-complex cone lemma. The exact partial theorem has passed an independent Xean review.
Same-model review, computation and compilation are not formal
verification or priority certification.

\subsection{Definitions, source checks, and actual intersection
algebra}\label{definitions-source-checks-and-actual-intersection-algebra}

The controlling book files are Chapter 72,
\texttt{ch72\_major\_open\_problems.tex} (polyhedral section, original
problem and residual scope), and Appendix M,
\texttt{appM\_ch18\_frontiers.tex} (analytic zero-current catalogue, no
product theorem used). Their complete paths and read-only-origin
snapshots are in \texttt{sources/} and \texttt{source-hashes.json}. The
exact original problem and completion criteria are preserved verbatim as
JSON strings in \texttt{task.json}, and the execution assignment in
\texttt{assignment.txt}. Neither book status nor prior review is a proof
input. Appendix M concerns the analytic continuous zero-current
frontier; no P06.2 result is used or claimed. BMPS, \emph{Monge--Ampère
measures on balanced polyhedral spaces}, arXiv:2603.08664v1, dated 10
March 2026, supplies definitions only. The explicit extension used here
has real function coefficients and real cycle weights; BMPS Definition
2.28 itself has integral slopes/rational constants. On any fixed finite
rational complex, affine-piece continuity is a rational linear system;
its real solution space has a rational basis, each made literal by
clearing slope denominators. Thus this convention is the real scalar
extension on that complex.

A PCC (BMPS Definitions 3.2--3.3) has rational endpoints
\(x_i\in X\cap N_{\mathbb{Q}}\), real coefficients \(\alpha_i\geq0\),
\(\sum\alpha_i=1\), and the actual, possibly irrational, centre
\(x=\sum\alpha_i x_i\). On some allowed rational structure there must be
faces \(\tau\preceq\sigma_i\) containing \(x\) and \(x_i\). A finite raw
PC function satisfies Jensen for \textbf{every} such combination, across
all allowed structures and remote endpoints. No continuity of raw
competitors is assumed.

With \(f=\gamma+u\), the raw envelope (Definition 6.2) is \[
 P_\gamma(u)(x)=\sup\{q(x):q\in\mathrm{PC}(X),\ q\leq f,
                          \ q\leq\gamma+C_q\text{ for some constant}\}.
                                                               \tag{4}
\] The last condition is only upper growth. The regularization class has
two-sided growth (Definitions 3.14, 3.23). Positive at infinity
(Definition 4.15) means strictly positive slope on every nonzero
recession vector of every affine reference cell. PAPC (Definition 3.1,
here with real weights) tests \(h\cdot z\geq0\) for every nonnegative
balanced real local cycle on every ordinary open, in every positive
dimension. On this two-dimensional carrier only dimensions one and two
can be nonzero.

Use the actual lattice
\(N_\sigma=N\cap\operatorname{span}_{\mathbb{R}}(\sigma-\sigma)\). For a
facet \(\tau\prec\sigma\), \(n_{\sigma/\tau}\) lifts the primitive
positive class in \(N_\sigma/N_\tau\). The divisor formula (Definition
2.31) is \[
 (h\cdot z)(\tau)=\sum_{\sigma\succ\tau}z(\sigma)h_\sigma(n_{\sigma/\tau})
       -h_\tau\left(\sum_{\sigma\succ\tau}z(\sigma)n_{\sigma/\tau}\right).
                                                               \tag{5}
\] The second sum is tangential by balancing. Changes of lift cancel
because the linear parts agree on tangents. Formula (5) is real
bilinear, depends only on incident germs, and annihilates ambient affine
functions.

Here is a direct reconstruction of refinement and balancing, rather than
an appeal to source intersection theory. For a new tested facet in a
refinement, take its minimal old face. If its dimension exceeds the
cycle dimension, inherited incident cycle weights are zero. If its
dimension equals the cycle dimension, the two inherited pieces in that
plane have equal weights and affine parts and opposite primitive
quotient normals, so cancel. If its dimension equals the tested facet
dimension, the actual tangent lattice and incident normal classes agree
with the old ones, so the old coefficient is inherited. These cases
prove refinement compatibility and inherited balancing. They apply
successively to both divisors; in particular new vertices in affine
edges cannot create a zero-cycle coefficient.

For balancing of the first divisor of a two-cycle, translate a vertex to
zero and take a pointed rational local fan, adding harmless rational
cuts if necessary. In a two-cone with primitive rays \(r,s\), let \[
 \chi_{rs}=[N_C:\mathbb{Z}r+\mathbb{Z}s],\qquad \omega_{rs}=z(C).
\] Extending primitive \(r\) to a basis of the saturated plane lattice
shows that its inward primitive quotient normal is represented by
\(s/\chi_{rs}\). A nonintegral representative has only a tangential
discrepancy from an integral lift, which cancels in (5). Top balancing
and the first weight are \[
 \sum_{C\succ r}(\omega_{rs}/\chi_{rs})s=t_r r,
 \qquad D_h(r)=\sum_{C\succ r}(\omega_{rs}/\chi_{rs})h(s)-t_rh(r). \tag{6}
\] Summing \(D_h(r)r\), the first terms, once per unordered cone, are
\(\omega_{rs}(h(s)r+h(r)s)/\chi_{rs}\). Substituting top balancing into
the subtracted terms gives exactly the same sum. Hence
\(\sum_rD_h(r)r=0\). This proves first-divisor balancing for real
weights in every germ. The full second coefficient is \[
                       (h^2\cdot[X])\{0\}=\sum_rD_h(r)h(r).     \tag{7}
\] Every intermediate ray, including zero-weight or zero-value rays,
belongs to this formula.

The same fan calculation proves \textbf{mixed commutativity}, rather
than importing BMPS Lemma 2.35's external intersection premise. At a
common structure vertex, write both affine germs as homogeneous
functions \(h,g\) after subtracting their values. Applying (5) to the
balanced one-cycle from (6) gives \[
 \begin{split}
 (g\cdot(h\cdot z))\{0\}
  &=\sum_{\{r,s\}}\frac{\omega_{rs}}{\chi_{rs}}
                 \bigl(h(s)g(r)+h(r)g(s)\bigr)
                   -\sum_r t_r h(r)g(r).
 \end{split}                                               \tag{M}
\] The right side is symmetric in \(h,g\), and includes all rays and
their actual saturated indices. For a one-cycle, its divisor is a
zero-cycle, whose balancing has no condition. Thus all recursive stages
needed on this two-dimensional carrier are balanced. Refinement
compatibility at each stage globalizes (M) and proves locality at every
zero-face. Adding an ambient affine function to either slot changes no
first divisor and gives zero against the balanced intermediate one-cycle
in the second slot, so also changes no mixed coefficient. This justifies
every affine subtraction in the contact proof. At a nonvertex, any
rational refinement that inserts a zero-face inherits zero second weight
unless the old structure already had that vertex. Consequently there is
no hidden measure on open lower faces or top-face interiors.

Product balancing itself holds in the actual lattice: at an edge times a
graph vertex, the quotient primitive normals are the primitive
directions of the second graph, and the weights are the first edge
weight times its balanced incident weights; the other type of facet is
symmetric. Interior subdivision facets have equal weights and opposite
normal classes. Thus the product is the actual balanced real two-cycle
used here. Definition 4.4's complete measure is the resulting finite
zero-cycle on structure vertices. On a unit complete plane, \(|s|+|t|\)
has first weight two on each axis ray, and (7) gives eight: \(2!\) times
the slope-square area four. No factorial division, probability
normalization, or singular-face Alexandrov substitution is used.

If a continuous real PA function is actual PC, it is PAPC against
\textbf{all} required cycles. Indeed, refine a local cycle and the
function, choose a rational point \(x\) in the ordinary open and
relative interior of its tested facet \(\tau\), and first take
nonnegative rational balanced incident weights \(c_\sigma\). Put
\(S=\sum c_\sigma>0\) and
\(T=\sum c_\sigma n_{\sigma/\tau}\in N_{\tau,\mathbb{Q}}\). For small
rational \(\varepsilon>0\), \[
 x_\sigma=x+\varepsilon(n_{\sigma/\tau}-T/S)
\] lies in its actual incident cell and open: the correction is
tangential and the quotient direction is inward. These rational
endpoints, with coefficients \(c_\sigma/S\), give an actual PCC centred
at \(x\). Jensen gives
\(0\leq\varepsilon(\sum c_\sigma h_\sigma(n_{\sigma/\tau})-h_\tau(T))\),
precisely (5). The cone of nonnegative balanced incident vectors is a
rational polyhedral cone, so has rational generators, as proved below.
Linearity extends the result to every real incident vector. The zero
vector gives zero. This proves all ordinary-open, all-dimension PAPC
tests. Balancing just proved and this positivity allow both recursive
divisors.

For all envelopes constructed here, \(P+1/n\) is decreasing uniform PAPC
approximation in \(\gamma+O(1)\). Constants have zero divisor; the
approximant measures are identical to \(P^2\cdot[X]\). Thus the
regularized operator is concretely the already defined PA operator,
without an approximation, gluing, locality or continuity theorem
imported from the source. The authentic proof bodies of BMPS Lemmas
2.34--2.35, Lemma 3.4, Proposition 3.22 and Proposition 6.3 were
inspected as comparison data. Their external locality/intersection
dependencies are not used. There are \textbf{no essential external
theorem premises} in the argument that follows.

\subsection{Elementary finite-dimensional facts used in the
reconstruction}\label{elementary-finite-dimensional-facts-used-in-the-reconstruction}

A polyhedral cone has finitely many generators; for rational
inequalities these can be rational. Remove its lineality (a rational
subspace for rational data) and work in the pointed quotient. The sum of
its nonnegative defining forms is strictly positive away from zero:
otherwise all forms vanish on a nonzero vector and its negative. The
section of value one is bounded, since an unbounded sequence, normalized
to the unit sphere, would yield a nonzero cone vector of value zero. A
bounded polyhedron is the convex hull of its vertices: in the minimal
face of any nonvertex, a line ends on smaller faces, and induction
reduces dimension. Its vertices are the finitely many full-rank
active-system solutions, rational for rational data. Lift back and add
both signs of a lineality basis. Homogenizing a nonempty pointed
rational polyhedron therefore proves \[
 \sigma=\operatorname{conv}\{v_1,\ldots,v_l\}
                    +\operatorname{cone}\{r_1,\ldots,r_m\},     \tag{8}
\] with rational vertices and recession vectors and at least one vertex:
the height-one slice needs a positive-height extreme ray, whose
extremality gives an original vertex. The homogenized cone is pointed
because any line has height zero and would give original lineality. The
same proof holds over real data without rational coordinates.

Eliminating a variable from finitely many nonstrict linear inequalities
compares every lower bound with every upper bound and retains
variable-independent conditions. A one-sided family imposes no further
conditions. Repetition proves that projections of polyhedra are
polyhedra; variable coefficients stay rational if originally rational,
even with real constants. It also proves closedness of finite
convex-hull-plus-cone sets by projecting their coefficient
representation.

Needed separations can be proved by nearest points. A point outside a
nonempty closed convex set has a nearest point after restriction to a
sufficiently large compact ball. Differentiating squared distance along
segments gives a separating normal. A compact convex set disjoint from a
closed convex set has positive attained distance, with the same compact
restriction for the second set; differentiating in both sets strictly
separates them. If the latter is a cone, a finite supremum of the
separating functional on it is zero. This proves the cone/polytope
separations used here.

If the pointed cone quotient is zero-dimensional, only its zero vector
and the lifted lineality basis are needed. The elementary fact about a
bounded polyhedron used above also covers its zero-dimensional case.
Pointedness ensures the homogenized cone has no lineality; its
zero-height extreme rays give recession generators and its
positive-height extreme rays give exactly the normalized vertices. This
proves (8) without an unstated decomposition theorem.

A finite rationally convex function on \(\mathbb{Q}\) has a unique
continuous convex extension on any interval interior. Rational endpoint
values bound it above on a rational compact interval. Reflection about a
rational centre bounds it below on a smaller interval. Bounds
\(L\leq F\leq M\) on a larger rational interval imply a Lipschitz bound
on a smaller one: extend the rational segment from \(y\) through \(z\)
by a rational factor \(R/|z-y|\) staying in the larger interval and use
Jensen, then reverse \(y,z\). Density gives the completion and
convexity. A finite convex function on \(\mathbb{R}\) is continuous by
this argument, and has a touching affine support at each point: its left
and right secant limits are finite, ordered, and any slope between them
gives the support.

For \(b(p)=\min_j(c_j+g_jp)\), the supergradient interval at \(p_0\) is
the convex hull of the active \(g_j\). Its directional derivative is the
minimum of their products with the direction. Separating a candidate
outside that interval disproves its supergradient property. Expanding
sums of finite minima proves that supergradients of nonnegative sums are
the corresponding weighted Minkowski sums, also when coefficients
vanish. The outward normal cone of a polytope is generated by active
inequality normals: an alleged normal outside their finitely generated
closed cone can be separated to give a direction satisfying the active
inequalities and increasing the alleged normal; a sufficiently small
step also satisfies inactive inequalities, a contradiction. These proofs
retain rational normal generators whenever inequality coefficients are
rational, even if constants are real.

\subsection{A finite rational cone for actual PC on a fixed
complex}\label{a-finite-rational-cone-for-actual-pc-on-a-fixed-complex}

This lemma concerns a \textbf{fixed} finite embedded rational complex
\(\Sigma\). It does not assert regularity on its carrier. Let
\(A(\Sigma)\) be the real vector space of continuous functions affine on
each cell, with real affine coefficients. Choose rational affine
coordinates on each cell's affine span. Agreement on each common face is
a finite rational homogeneous linear system: test agreement on a
rational affine basis of that face. These tests suffice because the
difference is affine. Thus \(A(\Sigma)\) has a rational basis, including
when cells have lineality or are lower dimensional. Clearing slope
denominators in a rational basis element puts it in the literal PA
class. This proves that the declared function space is its real scalar
extension.

For every nonempty face \(\tau\), every coface \(\sigma\succeq\tau\),
and \(y\in\operatorname{relint}(\tau)\), define \[
 T_{\tau\sigma}=\operatorname{cone}(\sigma-y).                 \tag{FC1}
\] Here cone means nonnegative scalar multiples; this set is already a
closed convex polyhedral cone. To prove this, write \(\sigma\) using
rational affine equalities and inequalities \(a_k(z)+b_k\geq0\). At
\(y\), an inequality is active if and only if it vanishes identically on
\(\tau\): a nonnegative affine function on \(\tau\) cannot vanish at a
relative-interior point unless its restriction is zero, by taking short
steps in both tangent directions. A vector \(d\) is in (FC1) exactly
when it satisfies the homogeneous equalities and \(a_k(d)\geq0\) for all
active inequalities. Necessity follows by substitution. For sufficiency,
every inactive inequality has positive slack at \(y\); there are
finitely many, so a sufficiently small positive step \(y+t d\) satisfies
all of them, as well as the active constraints. Therefore
\(y+t d\in\sigma\). This also proves that each such direction stays in
\(\sigma\) for every sufficiently small positive step. The same active
constraints hold at every real point in \(\operatorname{relint}(\tau)\),
proving independence of the chosen centre, rationality, and inclusion of
the whole tangent space of \(\tau\) as lineality. No boundedness,
pointedness, or full dimension was used.

Form the cone of vector families, indexed by \textbf{all} cofaces
including \(\tau\), \[
 B_\tau=\{(d_\sigma):d_\sigma\in T_{\tau\sigma},\quad
                              \sum_{\sigma\succeq\tau}d_\sigma=0\}.
                                                               \tag{FC2}
\] It is a rational polyhedral cone and has finitely many rational
generators, with both signs of a rational lineality basis retained. The
finite-generation proof above applies to the product of the tangent
cones and the rational sum-zero equations. Impose, for every generator
and every \(\tau\), \[
                \sum_{\sigma\succeq\tau}h_\sigma^{\rm lin}(d_\sigma)
                                  \geq0.                       \tag{FC3}
\] These are rational homogeneous linear inequalities in \(A(\Sigma)\).

\textbf{Claim.} Their solution cone \(C(\Sigma)\) equals
\(A(\Sigma)\cap\mathrm{PC}(X)\), with actual PC tested against all
allowed rational structures, all real centres, and remote rational
endpoints.

For necessity choose a rational \(y\in\operatorname{relint}(\tau)\).
Such a point exists: rational affine coordinates identify the relative
interior with a nonempty open subset cut out by rational constraints,
and rational points are dense in its affine span. Given a rational
generator of \(B_\tau\), let \(m\) be the number of cofaces. A
sufficiently small rational \(t>0\) puts all rational endpoints
\(y+mt d_\sigma\) in their cofaces. Their equal coefficients \(1/m\)
give a genuine \(\Sigma\)-PCC centred at \(y\). Since every piece agrees
at \(y\), Jensen is precisely (FC3) after division by \(t\). The zero
family causes no difficulty. Linear combination extends the inequalities
from generators to every real family in \(B_\tau\).

For sufficiency first take any straight segment contained in the
carrier, with endpoints allowed to be real. Each cell intersects it in
an interval, so a finite subdivision makes \(h\) affine on each
subinterval. At an interior breakpoint
\(x\in\operatorname{relint}(\tau)\), let \(d\) be the forward direction.
Its incoming \(-d\) and outgoing \(d\) lie in the tangent cones of their
respective initial cells, which contain \(\tau\) as a face. Put them in
the corresponding slots of (FC2). If the slots coincide, sum them in
that slot; both derivatives then come from the same affine piece and
agree. In all cases (FC3) gives \(h'_+(x;d)-h'_-(x;d)\geq0\). A
continuous PA function with nondecreasing successive slopes is convex:
the secant slope over an interval is a weighted average of its piece
slopes, so left secants do not exceed right secants. Thus \(h\) is
convex on every contained straight segment.

Now take any genuine PCC \((x,(x_i),(\alpha_i))\) on \textbf{another}
permitted structure. Each segment \([x,x_i]\) lies in its permitted
convex coface and hence in \(X\). If \(d_i=x_i-x\ne0\), its initial
short segment lies in some \(\Sigma\)-cell containing \(x\); finiteness
of the segment subdivision ensures such a cell exists. The minimal
\(\Sigma\)-face \(\tau\) of \(x\) is a face of that cell, and
\(d_i\in T_{\tau\sigma_i}\) by (FC1). Continuity and affine agreement
make the initial derivative unambiguous. Convexity along the segment
gives \(h(x_i)-h(x)\geq h'_{x}(d_i)\). Group
\(D_\sigma=\sum_{i:\sigma_i=\sigma}\alpha_i d_i\). Each group stays in
its convex tangent cone, and \(\sum D_\sigma=\sum\alpha_i(x_i-x)=0\).
Therefore (FC3) gives \[
 \sum_i\alpha_i(h(x_i)-h(x))\geq
 \sum_\sigma h_\sigma^{\rm lin}(D_\sigma)\geq0.              \tag{FC4}
\] Zero displacements and zero coefficients contribute zero. This proves
the claim at \textbf{every real centre}, with no insertion of irrational
endpoints or cycles into the source definitions.

\subsection{Exact positive-growth decomposition of the prescribed real
reference}\label{exact-positive-growth-decomposition-of-the-prescribed-real-reference}

Choose a common finite rational affine structure \(\Sigma\) for the
product and the prescribed real \(\gamma\). Since \(\gamma\) is positive
at infinity, each cell is pointed: nonzero lineality would put both
\(r\) and \(-r\) in its recession cone, requiring opposite affine slopes
both positive. By (8), each recession cone has finitely many nonzero
rational generators. Its positive-growth conditions are exactly the
finite strict inequalities \[
       h_\sigma^{\rm lin}(r)>0\quad
       \text{for all chosen nonzero recession generators}.     \tag{FC5}
\] A nonzero recession vector is a nonnegative combination with at least
one positive coefficient, so these inequalities suffice as well as being
necessary. Bounded cells impose none.

Write \(C(\Sigma)\) by its finite inequalities \(L_i\geq0\) in the
rational space \(A(\Sigma)\), and let \(I=\{i:L_i(\gamma)=0\}\),
\(S=\bigcap_{i\in I}\ker L_i\). Then \(F=C(\Sigma)\cap S\) is the
minimal face containing \(\gamma\). Indeed it is a face by nonnegativity
of the \(L_i\). In \(S\), all other inequalities are strictly positive
at \(\gamma\), so a small relative ball about \(\gamma\) lies in \(F\).
Hence \(\operatorname{span}F=S\) and
\(\gamma\in\operatorname{relint}F\). Any face containing \(\gamma\)
contains that ball: for \(z\in F\) a short step from \(\gamma\) opposite
\(z-\gamma\) stays in the ball, expressing \(\gamma\) as a strict convex
combination of \(z\) and another point of \(F\); the defining face
property then includes \(z\). This proves minimality without a face
theorem. The subspace \(S\) is rational, even when the cone has
lineality.

Choose a rational basis of \(S\), and use it as coordinates. If its
dimension is \(d>0\), choose rational endpoints \(a_k<t_k<b_k\) about
the real coordinates \(t_k\) of \(\gamma\), so close that their entire
box lies in the preceding relative ball and satisfies every strict
inequality (FC5). Density of the rationals and finiteness of all
constraints give such a box. All its \(2^d\) vertices \(\gamma_j\) have
rational affine coefficients, belong to \(C(\Sigma)\), and retain every
strict recession slope. Put \(\theta_k=(t_k-a_k)/(b_k-a_k)\in(0,1)\).
The product weights obtained by choosing \(\theta_k\) at an upper
endpoint and \(1-\theta_k\) at a lower endpoint are strictly positive
and sum to one. Their coordinate averages are exactly \(t_k\). Thus,
\textbf{as functions on the entire unchanged carrier}, \[
                 \gamma=\sum_{j=1}^{J}\lambda_j\gamma_j,
                 \lambda_j>0,\qquad\sum_j\lambda_j=1.        \tag{FC6}
\] Each \(\gamma_j\) is actual PC against every structure by (FC4), not
just feasible in a nodal relaxation. If \(d=0\), then \(\gamma=0\) in
this space; use the single rational reference \(0\) with weight one. In
that case (FC5), if nonempty, would exclude \(0\), so this case can
occur only when there are no nonzero recession directions. This handles
the degenerate case without claiming growth where it is impossible. No
approximation, reference replacement, or lost strict recession
inequality occurs in (FC6).

\subsection{Graph geometry, real weights, and genuine
cofaces}\label{graph-geometry-real-weights-and-genuine-cofaces}

At a vertex the full kernel of the rational direction matrix is
one-dimensional. Otherwise perturb its strictly positive actual weight
vector by a sufficiently small nonproportional kernel vector; positivity
persists, contradicting (1). The kernel therefore has a primitive
positive integer generator \(\lambda_v\), and the actual incident vector
is \(a_v\lambda_v\), \(a_v>0\). At a bivalent vertex the directions are
opposite primitive vectors and weights equal. At a branched vertex there
is no opposite pair, since that pair would give another nonnegative
relation. No one-valent vertex can balance.

On a connected graph all edge weights are a single positive real scalar
times positive rational weights. Fix one edge weight. At each adjacent
vertex the ratios of incident weights equal ratios of positive integers
in \(\lambda_v\); propagate along paths through the finite graph.
Existing actual weights ensure consistency around every cycle. This
yields rational normalized weights \(r_e\) with \(w_e=c r_e\). Full-line
components have one arbitrary positive scalar weight. Rational
subdivision preserves these statements. This reduction is essential: it
makes the incidence elimination below have rational variable
coefficients, without supposing the given real scale rational.

Every connected unbranched component is a complete rational affine line.
Propagating through bivalent vertices keeps every edge on the same
affine line with equal weight. The component is a connected finite
closed union of intervals/rays on that line. Any finite endpoint of its
union would be one-valent, contrary to balancing. It therefore covers
the whole line; a bounded cycle cannot be embedded into this line as a
pure polyhedral graph. A component supplied already as a full line has
the same conclusion. Its tangent lattice is the saturated primitive line
lattice \(\mathbb{Z}v\), independently of the ambient lattice's
complementary directions.

For the all-coface PCC tests on a cylinder \(G\times\mathbb{R}\), any
convex cell projects to a convex subset of the graph. A segment cannot
pass a branched vertex since it would need an opposite pair there. Thus
its graph projection is in a single straight chain through bivalent
subdivisions or in a vertex. If the centre's graph coordinate is in a
straight-chain interior, all cofaces containing the common face remain
in that chain and the graph test is ordinary one-variable convexity. At
a branched centre \(v\), each coface projects into one incident straight
chain; every endpoint displacement has the form \(t_i v_e\),
\(t_i\geq0\), up to the next branched vertex. Grouping the barycentric
displacements gives \[
 b_e=\sum_{i\text{ in }e}\alpha_i t_i,
 \qquad\sum_e b_ev_e=0,\qquad b_e=c r_e,\ c\geq0.           \tag{9}
\] This argument uses actual cofaces on arbitrary rational structures,
not merely a prescribed product subdivision. Distant endpoints and
inserted bivalent vertices are included. Zero graph displacements cause
no contribution.

\subsection{Positive direction: full graph-cylinder reconstruction in
dimension
two}\label{positive-direction-full-graph-cylinder-reconstruction-in-dimension-two}

Suppose one factor is unbranched. Its components are complete rational
affine lines of constant positive real weights. Work first on one
connected component \(G\times\mathbb{R}_z\). Use its actual primitive
coordinate \(z\), and normalize graph weights to positive rational
\(r_e\); the actual product two-cycle is a positive scalar multiple of
these normalized weights. This scalar affects neither raw PC nor contact
and will be restored in the measure. No hypothesis of separability or
ambient convex extension of the data is made.

\textbf{Scalar graph lemma.} A finite function on rational graph points
satisfying rational-coefficient graph PCCs has a continuous
edgewise-convex completion \(k\) with finite outgoing derivatives and \[
                              \sum_e r_e k'_e(v+)\geq0.        \tag{10}
\] Complete it on edge interiors by the rational convex fact. The
one-sided limit at an endpoint is finite and at most its specified
endpoint value: endpoint secants give the upper bound and a fixed
interior secant gives a lower bound. For small rational \(t>0\), the
points \(v+t v_e\), with rational coefficients \(r_e/\sum r_e\), form an
actual vertex PCC. Its inequality and those endpoint upper bounds force
every limit to equal \(k(v)\). Each outgoing derivative is bounded above
by a finite secant. Dividing the same PCC by \(t\) prevents any
derivative from being \(-\infty\), and yields (10).

Conversely, any continuous edgewise-convex \(k\) satisfying (10), plus
\(pz\) for any real \(p\), satisfies \textbf{every} actual cylinder PCC.
On a chain this is convexity (the bivalent inequality joins the
derivatives). At a branched centre use (9) and
\(k(q_i)-k(v)\geq t_i k'_e(v+)\); summation gives
\(c\sum r_e k'_e(v+)\geq0\). Affine \(z\) terms cancel by the barycentre
identity. This includes arbitrary real centres/coefficients and every
coface structure.

For an arbitrary finite raw cylinder PC function \(h\), define on
rational graph points \[
 k_{h,p}(q)=\inf_{z\in\mathbb{Q}}(h(q,z)-pz).                 \tag{11}
\] It satisfies rational graph PCCs: choose rational \(z_i\) tending to
each endpoint infimum, apply the product PCC with rational coefficients
at \((q,\sum\alpha_i z_i)\), subtract its affine term, and take the
infimum. A \(-\infty\) endpoint is handled by varying that endpoint
alone. If (11) is finite at one rational graph point, it is finite
throughout the component. From an edge-interior finite centre any other
rational edge point can be a positively weighted endpoint, with the
other endpoint on the opposite side. From a finite vertex, an arbitrary
rational point \(v+t_ev_e\) on an incident edge is an endpoint in a
positive rational vertex PCC: choose positive rational distances on the
other edges and coefficients proportional to \(r_e/t_e\). No endpoint of
a finite centre's such inequality can be \(-\infty\). Propagation across
the connected graph includes every cycle and end. Thus the scalar lemma
applies whenever (11) is finite anywhere.

Put \(f=\gamma+u\), \(C=\|u\|_\infty\), and take a common finite
rational spatial structure for \(\gamma,f\), refined by edge strips.
Split full-line edge components at a rational point if needed. Every
cell is pointed: lineality would give opposite nonzero reference
recession vectors with opposite affine slopes, both supposedly positive.
Define \[
 D=\{p\in\mathbb{R}:p r_z\leq s_r\text{ for every cell recession generator}
                                      \ r\text{ with }r_G=0\},           \tag{12}
\] where \(s_r\) is the affine obstacle slope. Bounded \(u\) has zero
recession slope, so these are the same constraints for \(\gamma\). They
are finite, have rational variable coefficients and real constants, and
give a closed interval with zero in its interior by strict growth. For
\(p\in D\), \[
                         F_p(q)=\inf_{z\in\mathbb{R}}(f(q,z)-pz)          \tag{13}
\] is finite and attained at a vertex of some spatial cell fibre. The
fibre is pointed. Its recession cone is generated by zero-graph cell
rays: on a bounded strip all graph recession is zero; on an end strip it
is nonnegative, so positive graph rays cannot combine to yield zero
graph recession. The objective is nonnegative on the fibre recession
cone, so (8) supplies a minimizing vertex. There are finitely many
cells. If \(p\notin D\), a violating rational ray based at a rational
cell point gives \(F_p(q)=-\infty\) at a rational graph point. The
interval \(D\) is bounded: at any fixed rational graph point, finitely
many pieces give \(f(q,z)\leq M|z|+B\); if \(|p|>M\), (13) is
\(-\infty\). Therefore \(D\) is compact.

Subdivide graph edges at the projections of all spatial vertices. These
projections are rational and finite. Between consecutive projections, a
fibre vertex lies on a spatial one-face: intersection with the slicing
hyperplane lowers its minimal-face dimension by at most one, so a fibre
vertex's minimal face has dimension at most one. Such a traversing face
has rational trace \(z(t)=z_i+t b_i\); its affine obstacle height has
real coefficients. Its feasibility throughout that interval is constant,
since changes occur only at projected face vertices. Consequently \[
 F_p(t)=\min_i(A_i(p)+tB_i(p)),                              \tag{14}
\] with finitely many indices independent of \(p\), and affine
\(A_i,B_i\) with rational \(p\)-gradients and real constants. At a graph
node, direct fibre slicing gives the same kind of finite affine minimum.
Formula (14) extends continuously and concavely to the finite endpoints:
finite candidate limits have objectives at least \(F_p(v)\); conversely,
fixing any finite \(z\) and using continuity of \(f\) gives
\(\limsup F_p(t)\leq f(v,z)-pz\), then take its infimum. Thus changes of
endpoint fibre type are retained.

After the last projected vertex on an end, take its primitive parameter
\(t\geq0\) from that node. Define \[
 \theta_e(p)=\min_iB_i(p)=\lim_{t\to\infty}F_p(t)/t.          \tag{15}
\] Every \(B_i-\theta_e\geq0\), so \(F_p(t)-t\theta_e\) is
nondecreasing. Hence the line with slope \(\theta_e\) starting below
\(F_p(0)\) lies below the entire tail obstacle. All tails and pure/mixed
recession constraints are included.

For fixed \(p\in D\), on the finite graph node set \(V\) impose the
incidence system \[
 a_v\leq F_p(v),\qquad
 L_v(a,p)=\sum_{e=vw\text{ bounded}}\frac{r_e}{\ell_e}(a_w-a_v)
              +\sum_{e\text{ tail at }v}r_e\theta_e(p)\geq0.             \tag{16}
\] The actual primitive lengths \(\ell_e\) are positive rational
numbers. Every bounded cycle edge appears in this system; no graph is
unfolded.

It is feasible. The global competitor \(\gamma-C\leq f\) is actual PC.
Its partial infimum is finite on \(D\) by (12)--(13), and the scalar
lemma applies. Replace it on every bounded interval by its endpoint
chord and on every tail by its endpoint line of slope \(\theta_e\).
Convexity makes these replacements dominate it. Concavity of (14) keeps
the chords below \(F_p\); (15) keeps the tail lines below it. At both
endpoints the new outgoing chord derivative is at least the old
derivative. Convexity and the tail asymptotic upper bound force all old
tail secants to be at most \(\theta_e\), so its outgoing derivative also
increases. This proves (16).

Coordinatewise maxima preserve feasibility: at a node chosen from one
vector all neighbour values only increase, so its \(L_v\) does not
decrease. Choose a feasible witness \(b\); maximize \(\sum a_v\) in the
compact box \(b_v\leq a_v\leq F_p(v)\). If any feasible vector exceeds a
maximizing vector in one coordinate, their coordinatewise maximum is
feasible in that box and increases the objective. There is therefore a
greatest feasible vector \(a(p)\). Its chord/tail interpolation
\(k_p(q)\) is continuous, obeys (10), lies below \(F_p\), and dominates
every scalar function (11) lying below \(F_p\), by the same chord/tail
replacements.

Each \(a_v\) is a finite concave PA function on \(D\) with rational
gradients. Expand all caps into their affine minima and each weighted
sum of tail minima into a minimum over all choices of pieces. This turns
(16) into finitely many linear inequalities in \((a,p)\) with rational
variable coefficients and real constants. Eliminate all \(a_w\),
\(w\ne v\), by the proved elimination procedure. For fixed \(p\), the
section is \((-\infty,a_v(p)]\), since uniformly lowering all node
values preserves the Laplacians. Thus all nontrivial bounds on \(a_v\)
are upper bounds; zero-coefficient inequalities only restrict \(D\).
Dividing positive rational coefficients yields \[
                       a_v(p)=\min_j(c_{v,j}+g_{v,j}p),\quad
                       g_{v,j}\in\mathbb{Q}.                 \tag{17}
\] This proves continuity through the boundary of \(D\), as well as
rational gradients. The same holds for \(\theta_e\).

Set \[
                            Q(q,z)=\sup_{p\in D}(k_p(q)+pz).             \tag{18}
\] It is finite by compactness and continuity and actual PC by the
scalar lemma and the supremum Jensen inequality. It is at most \(f\).
For any raw \(h\leq f\) at rational \((q,z)\), the rational slice
\(h(q,\cdot)\) has the continuous convex completion proved above. A
touching affine support of slope \(p\) gives \(h(q,z)=k_{h,p}(q)+pz\).
The partial infimum is finite there, hence throughout the component by
propagation. A violating recession ray would contradict \(h\leq f\), so
\(p\in D\). Greatest-assignment domination gives \(h(q,z)\leq Q(q,z)\).

Partition \(D\) at the finitely many breakpoints of (17) and (15); let
\(J\) be the finite set of chamber endpoints. They may be real. On any
graph interval, \(k_p(q)+pz\) is affine in \(p\) on each chamber;
therefore (18) is the maximum over \(p\in J\) of finitely many real
affine strip functions. These agree at graph nodes. Hence \(Q\) is
continuous. At any real carrier point choose its minimal face on a
rational structure; unless it is a rational zero-face, rational vertices
of arbitrarily small enclosing simplices in that face have the point as
a real barycentre. The same-cell PCC and continuity of \(Q\) extend
\(h\leq Q\) to that point, without raw continuity. Applying this to
\(h=\gamma-C\) gives \[
                         \gamma-C\leq Q\leq f\leq\gamma+C.              \tag{19}
\] Thus \(Q\) itself is a raw competitor and bounds all competitors,
proving \(P_\gamma(u)=Q\) exactly.

\textbf{Rational spatial domains.} A finite real affine maximum alone
does not justify rational walls. Fix \(p_0\in J\). Let \(A_v\) be the
supergradient interval of \(a_v\) at \(p_0\), \(B_e\) that of
\(\theta_e\), and \(N_D\) the outward normal cone of \(D\) there.
Endpoints of \(A_v,B_e\) are rational by (17), and \(N_D\) has rational
generators by (12). For finite concave polyhedral \(b\), \(p_0\)
maximizes \(b(p)+pz\) on \(D\) exactly when \[
                                z\in-\partial^+b(p_0)+N_D.             \tag{20}
\] Sufficiency is the supergradient/normal inequality. If necessity
failed, separating the compact interval \(z+\partial^+b\) from the
closed cone \(N_D\) gives a direction nonpositive on all active normals
and strictly positive on every objective directional term. A small
feasible step raises the objective, a contradiction. This proves the
boundary cases too.

On a bounded edge \(0\leq t\leq\ell\), its full maximization domain is
\[
 z\in-(1-t/\ell)A_v-(t/\ell)A_w+N_D.                       \tag{21}
\] It is the convex hull of rational points
\((0,-\operatorname{Vert}A_v)\), \((\ell,-\operatorname{Vert}A_w)\),
plus the cone of rational vertical normal generators. On a tail the
domain is \[
 t\geq0,\qquad z\in-A_v-tB_e+N_D,                         \tag{22}
\] generated by rational base points and rational rays
\((1,-\operatorname{Vert}B_e)\) and vertical normal rays. Expanding the
finite weighted minima proves these supergradient formulas even at zero
coefficients. The domains are finite rational spatial polyhedra. Maxima
domains meet in equality faces; including their faces and removing
duplicates gives a rational complex. At a graph node every incident
formula has the same fibre domain \(-A_v+N_D\), so complexes glue.
Generated bounded-strip domains have vertices only at graph endpoints,
and tail domains only at their base. Thus \textbf{every} possible
zero-cycle vertex lies over a node of \(V\); there are no omitted
open-edge or distant-tail vertices.

By the PC-to-PAPC argument following (7), \(P\) is PAPC against all
local real cycles. The full recursive measure is finite and nonnegative.
The sequence \(P+1/n\) proves actual regularizability with exactly this
measure.

\textbf{Complete contact.} At any possible atom \(x=(v,z_0)\), suppose
\(P(x)<f(x)\). Define the nonempty compact exposed interval \[
 C_x=\{p\in D:a_v(p)+pz_0=P(v,z_0)\}.                      \tag{23}
\] It is convex because \(a_v\) is concave and the displayed expression
never exceeds its maximum. On \(C_x\), \(a_v=P(x)-pz_0\) is affine. Put
\[
 s_e(p)=(a_w(p)-a_v(p))/\ell_e\quad\text{on bounded outgoing edges},
 \qquad s_e(p)=\theta_e(p)\quad\text{on tails}.              \tag{24}
\] These are continuous concave polyhedral functions on \(C_x\), and
\(\sum r_es_e\geq0\). In local outgoing coordinates \(t\geq0\),
\(\zeta=z-z_0\), the homogeneous germ is \[
                            h_e(t,\zeta)=\max_{p\in C_x}(p\zeta+ts_e(p)).\tag{25}
\] Indeed the finite set \(J\) expresses all strip functions. Inactive
planes at \(x\) have positive gaps and stay inactive locally. Every
\(p\in C_x\) is a convex combination of endpoints of a chamber on which
all terms are affine; equality at the maximum forces its relevant
endpoints to be active. This proves the exact germ identity.

If \(C_x\) is a single point \(p_0\), after subtracting \(p_0\zeta\) the
germ is independent of \(\zeta\) on every branch. Its first divisor is a
constant-weight vertical line over the graph vertex (all interior branch
planes are affine); the function restricts to an affine function on that
line. Formula (5) gives zero second coefficient, or equivalently its two
opposite primitive vertical rays in (7) cancel. This retains artificial
cuts and all first faces by refinement compatibility.

If \(C_x\) is a nondegenerate interval, for every interior \(p\) choose
\(\eta>0\) with \([p-\eta,p+\eta]\subset C_x\). Globally in the vertex
fibre, \[
 P(v,y)-py\geq P(v,z_0)-pz_0+\eta|y-z_0|.                 \tag{26}
\] Thus continuous \(f(v,y)-py\) is coercive and attains its minimum. If
\(a_v(p)=F_p(v)\), (26) forces its minimizing point to be \(z_0\) and
equality \(f(x)=P(x)\), a contradiction. Hence \(a_v(p)<F_p(v)\)
throughout the interval interior. Greatest-assignment maximality forces
\(L_v(a(p),p)=0\): if positive, increase only \(a_v\) by a small amount
preserving its strict cap and that Laplacian. Every neighbour Laplacian
increases, tails are unchanged, and feasibility persists, contradicting
maximality. If there are no bounded incident edges, such an increase
changes no Laplacian and already contradicts the strict cap. Continuity
now gives \(\sum r_es_e=0\) on all of \(C_x\).

Positive weights and concavity force every \(s_e\) affine on this
interval: their nonnegative Jensen deficits have weighted sum zero. Its
affine gradient is rational, because an interior subinterval avoids the
finitely many walls of (17) and (15), whose gradients and primitive
lengths are rational. Write \(s_e(p)=\alpha_e+\beta_ep\),
\(\beta_e\in\mathbb{Q}\), \(\alpha_e\in\mathbb{R}\). The zero sum gives
\(\sum r_e\alpha_e=\sum r_e\beta_e=0\). Since the star direction kernel
is exactly this one relation, there exist a real linear form \(A\) and a
rational linear form \(B\) on its actual rational span, with
\(A(v_e)=\alpha_e\), \(B(v_e)=\beta_e\). This follows by factoring the
prescribed values through the rational direction matrix, not by
substituting a ray-generated lattice.

The full star germ is therefore \[
                              A(s)+\operatorname{supp}_{C_x}(\zeta+B(s)).\tag{27}
\] Extend these forms to the ambient space, rationally for \(B\). After
removing \(A\) and an affine multiple of the rational form
\(\ell=\zeta+B(s)\), the residual is a one-variable convex homogeneous
function, affine on \(\ell\geq0\) and \(\ell\leq0\), and zero on
\(\ell=0\). Refine the actual fan by this rational hyperplane. Its first
divisor vanishes away from \(\ell=0\), since the residual is ambient
affine there. On the entire support of the first divisor the residual is
zero, so its second divisor is zero by (5). This is the complete
recursion, including singular axes and all indices; it is not a
top-face-interior comparison. Consequently every noncontact atom has
coefficient zero. Finite support and bounded gap (19) prove (2).

For finitely many components construct the data, envelopes and measures
separately. A PCC's cells all contain its common face, so lie in one
connected component; a local cycle can similarly be treated
componentwise. Components are disjoint finite closed polyhedral subsets
and are open relative to their union. The finite component constructions
give a finite rational structure, a maximum of their finitely many
growth constants, and measures whose contact integrals add. Restore the
actual positive real product scalar on each component; zero stays zero.
This finishes the positive direction for all stipulated data. This proof
was reconstructed here in full; its predecessor's accepted review is not
invoked.

\subsection{Auxiliary rational-reference localization: whole-space
growth}\label{auxiliary-rational-reference-localization-whole-space-growth}

In this auxiliary construction only, write \(\eta\) for any
rational-coefficient positive-growth actual-PC reference. Its elementary
construction will be applied separately to every component \(\gamma_j\)
of (FC6); it is not a replacement of the prescribed real reference.

Assume both graphs branch, and choose original branched vertices
\(A,B\). All objects in this construction remain on the \textbf{entire
original product}. Choose ambient lattice bases and put \[
\rho(p,q)=\|p-A\|_1+\|q-B\|_1,\qquad
\rho_0(r)=\|r_1\|_1+\|r_2\|_1.                               \tag{28}
\] The second is the \textbf{homogeneous recession norm}, not the
translated function \(\rho\). Refine the original product and \(\eta\)
into a common finite rational structure \(\Sigma\), splitting any
full-line carrier pieces if needed. Each nonempty cell is pointed.
Indeed a lineality vector \(r\ne0\) and \(-r\) would both be recession
vectors of an affine reference cell, and its affine slopes on them are
opposite; they cannot both be strictly positive. Any subcell inherits
positive slopes on its nonzero recession vectors.

By (8), write every cell as the convex hull of finitely many rational
vertices \(v_i\) plus the cone of finitely many nonzero rational
recession generators \(r_j\). Its affine slope
\(\eta_\sigma^{\mathrm{lin}}(r_j)\) is a positive rational number.
Select one rational \(m>0\) no larger than any of the finitely many
ratios \[
       \eta_\sigma^{\mathrm{lin}}(r_j)/\rho_0(r_j).            \tag{29}
\] If there are no recession generators, choose any rational \(m>0\).
For \(x=\sum_i\beta_iv_i+\sum_jt_jr_j\), \(\beta_i,t_j\geq0\),
\(\sum_i\beta_i=1\), affinity and the norm triangle inequality give \[
\begin{split}
\eta(x)&\geq\min_i\eta(v_i)+m\sum_jt_j\rho_0(r_j),\\
\rho(x)&\leq\max_i\rho(v_i)+\sum_jt_j\rho_0(r_j).
\end{split}                                                        \tag{30}
\] Choose rational \[
B=\max\bigl(0,\ \max_{\sigma\in\Sigma}
                 [m\max_{v\in\operatorname{Vert}\sigma}\rho(v)
                        -\min_{v\in\operatorname{Vert}\sigma}\eta(v)]\bigr).
\] The finite vertex sets are nonempty, and all values in this
expression are rational. Equations (30) prove on \textbf{all}
components, bounded cells and ends that \[
                         \eta\geq m\rho-B.                 \tag{31}
\] This derives the needed coercivity from the exact recession
hypothesis, rather than assuming a global growth theorem or using a
conveniently chosen reference.

\subsection{Rational maps on the original stars and a potential in the
fixed
class}\label{rational-maps-on-the-original-stars-and-a-potential-in-the-fixed-class}

By the graph kernel proof, the primitive directions at \(A,B\) have sole
primitive positive integer relations \(\lambda_i,\mu_j\), and actual
incident weights \(a\lambda_i,b\mu_j\), \(a,b>0\). Partition each
incident index set into three nonempty groups \(I_k,J_k\), \(k=1,2,3\).
Set \[
 S_k=\sum_{i\in I_k}\lambda_i,\quad T_k=\sum_{j\in J_k}\mu_j,
 \quad\nu_1=(1,0),\ \nu_2=(0,1),\ \nu_3=(-1,-1),\quad c_k=1,1,2.
\] Prescribe on the original directions \[
L_1(v_i)=\nu_k/S_k\quad(i\in I_k),\qquad
L_2(w_j)=\nu_k/T_k\quad(j\in J_k).                            \tag{32}
\] Their weighted sums are \(\sum_k\nu_k=0\). Since the \textbf{entire}
rational direction kernel is the single relation, these assignments
annihilate it, factor through the actual rational spans, and extend to
rational ambient linear maps by rational basis extension. This is linear
algebra on the original spans, not a lattice or cycle substitution.

Define the global ambient convex rational-PA function \[
             H(p,q)=\|L_1(p-A)-L_2(q-B)\|_1.                 \tag{33}
\] For a matrix \(T\), the triangle inequality proves
\(\|Tx\|_1\leq(\max_j\sum_i|T_{ij}|)\|x\|_1\). Choose positive rational
\(L\) at least the two indicated matrix norms; then \(H\leq L\rho\).
Select rational \(0<\varepsilon<m/L\), and rational
\(C>\max\{0,\eta(A,B)+2\}\). Put \[
                  P=\max\{\eta,\varepsilon H+C\}.          \tag{34}
\] Ambient affine functions have equality in every admitted PCC, and
their finite maximum gives Jensen; thus \(H\) is actual PC. A maximum of
two PC functions is PC: at any centre select an active function, apply
its Jensen inequality, and bound each endpoint value by the maximum.
Hence \(P\) satisfies \textbf{every genuine PCC}, not just coarse
product or nodal tests. Both functions in (34) have rational affine
coefficients. On a common rational structure their comparisons are
rational affine inequalities. Subdividing by these inequalities and the
finitely many rational affine-piece walls, then including faces, gives a
finite rational structure \(\Pi\) on which \(P,\eta,H\) are affine. The
same observation proves all rational-domain assertions below. The same
assertion about comparisons would fail for arbitrary real coefficients;
(FC6) and the positive sum below avoid that assertion.

By (31), \[
 \eta-(\varepsilon H+C)\geq(m-\varepsilon L)\rho-B-C.
\] Take a rational \(R>(B+C)/(m-\varepsilon L)\). Then \(P=\eta\) on
\(\rho>R\). Consequently \(P-\eta\geq0\) is bounded and compactly
supported in the closed finite polyhedral carrier intersected with
\(\rho\leq R\), a compact Euclidean set. It follows that
\(P\in\eta+O(1)\). On every common affine cell and recession vector, the
affine difference \(P-\eta\) has slope zero, since it is globally
bounded; therefore \(P\) retains positive slopes at infinity. Its cells
are pointed by the same opposite-vector argument. The PC-to-PAPC and
two-stage reconstruction above show that \(P\) is actual
\(\mathrm{PAPC}_{\mathbb{R}}\), with finite nonnegative PA operator, and
\(P+1/n\) is decreasing uniform PAPC regularization in the fixed
\(\eta+O(1)\) class, with exactly the same measures. This proves the
operator domain without an unavailable approximation theorem.

\subsection{Positive sums of rational
localizations}\label{positive-sums-of-rational-localizations}

We now apply the auxiliary construction (28)--(34) to \textbf{each}
rational reference \(\eta=\gamma_j\) in (FC6), using the \textbf{same}
selected original branched junctions \(A,B\), groups and rational maps
defining \(H\). For each \(j\), choose its rational
\(m_j,B_j,\varepsilon_j,C_j\) as in (29)--(34), and obtain \[
 P_j=\max\{\gamma_j,\varepsilon_j H+C_j\},\quad
 P=\sum_j\lambda_jP_j,\quad
 C=\sum_j\lambda_jC_j,\quad
 \varepsilon=\sum_j\lambda_j\varepsilon_j>0.                \tag{RS1}
\] Each \(P_j\) is actual PC and rational PA on a finite rational
structure, dominates \(\gamma_j\), and equals \(\gamma_j\) off a compact
set. Positive sums preserve every Jensen inequality. Taking a common
rational refinement of the finitely many \(P_j\) structures, the
\textbf{sum} \(P\) is continuous real PA on finite rational domains and
actual PC. We never take \(\max\{\gamma,\varepsilon H+C\}\), nor infer
rational walls from the coefficients of a real maximum. Equations (FC6)
and (RS1) prove \(P\geq\gamma\), \(P=\gamma\) outside the finite union
of those compact sets, and \(P\in\gamma+O(1)\). On a common affine cell,
boundedness of \(P-\gamma\) makes its slope zero on every recession
vector, so \(P\) retains the prescribed strictly positive recession
slopes. Its cells are pointed. The all-local-cycle PC-to-PAPC proof and
explicit \(P+1/n\) regularization apply without any arithmetic
restriction on \(\lambda_j\).

\subsection{A common original product junction and the compact
obstacle}\label{a-common-original-product-junction-and-the-compact-obstacle}

For every \(j\), the choice \(C_j>\gamma_j(o)+2\), \(o=(A,B)\), and
continuity give \(\gamma_j-C_j<0\) near \(o\). In that neighbourhood
\(P_j=C_j+\varepsilon_jH\), hence \(P=C+\varepsilon H\). Choose the
common rational structure \(\Pi\) by cutting the original product with a
finite global collection of rational affine hyperplanes extending the
fixed spatial domains of all \(\gamma_j\), the \(H\) walls, and their
rational comparison walls. This procedure refines each function
structure: cut by all supporting hyperplanes of its cell boundaries; a
new open cell cannot cross any such boundary, so is contained in an old
cell. Intersections and faces remain finite and rational. It introduces
no wall depending on the real \(\lambda_j\).

Keep \(o\) as a vertex of the original product. Choose rational
\(r_0>0\) small enough for all \(j\) inequalities above, for the carrier
inside \(\rho<r_0\) to be its original short product star, and to avoid
all hyperplanes not passing through \(o\), all nonincident cells and all
incident bounded endpoints. These finitely many closed sets not
containing \(o\) have positive distances from it; take the minimum with
the finitely many continuity radii. Within this ball every active
spatial constraint is homogeneous about \(o\). At any \(x\ne o\), small
expansion and contraction along the radial line stay in its minimal
face: homogeneous equalities remain equalities and the remaining strict
inequalities stay strict. Thus \(x\) is not a vertex, and \[
 P=C+\varepsilon H\quad(\rho<r_0),\qquad
                  \Pi^{(0)}\cap\{\rho<r_0\}=\{o\}.           \tag{35}
\] The ball is always intersected with the original carrier. This
argument does not equate an ambient ball with a product star.

Choose a sufficiently small rational \(\tau>0\) so that \[
 z_{ij}=(A+\tau S_kv_i,\ B+\tau T_kw_j),\qquad
 \alpha_{ij}=\frac{\lambda_i\mu_j}{3S_kT_k},
 \quad i\in I_k,\ j\in J_k                                  \tag{36}
\] lie in their original product cofaces, before bounded endpoints, and
in \(\rho<r_0\). In (36), \(\lambda_i,\mu_j\) are the \textbf{integer
star relations}, distinct from the real convex-decomposition weights in
(RS1). Each group has total coefficient \(1/3\); the sum is one. The
first weighted displacement is \(\tau\sum_i\lambda_iv_i/3=0\), and the
second is \(\tau\sum_j\mu_jw_j/3=0\). Thus this is a genuine
original-product PCC with common face \(o\). Both maps at an endpoint
give \(\tau\nu_k\), so \(H(z_{ij})=0\) and \(P(z_{ij})=C=P(o)\). The
rational endpoints need not be in a common fine \(P\)-coface, since raw
PC also tests the original rational structure.

Choose integer \(M\) with \(1/M<r_0\) and \(M\rho(z_{ij})\geq1\) at
every nonzero endpoint. Put \[
             b_0=(1-M\rho)_+,\qquad u=P-\gamma+b_0.           \tag{37}
\] The norm has rational spatial walls and rational coefficients, and
cutting by \(\rho=1/M\) on its finitely many pieces remains rational.
Therefore \(u\) is nonnegative, bounded, continuous real PA on a finite
rational structure, with compact support. Its ceiling for the
\textbf{unchanged original real reference} is \(\gamma+u=P+b_0\). The
bump is one at \(o\), zero at every endpoint (36), and zero at every
other \(\Pi\)-vertex. The reference and potential agree at infinity on
every component.

\subsection{Exact raw envelope on all components, cycles and
ends}\label{exact-raw-envelope-on-all-components-cycles-and-ends}

For any raw competitor \(q\leq\gamma+u=P+b_0\), the endpoint values in
(36) are at most \(C\). Their genuine PCC gives \(q(o)\leq C=P(o)\). At
every other vertex of \(\Pi\), the bump is zero, giving
\(q(v)\leq P(v)\). The \textbf{sole upper growth} condition is
\(q\leq\gamma+C_q\leq P+C_q\). No continuity, lower bound, local gluing
or graph unfolding is imposed on \(q\).

Let \(x\) be an arbitrary real point of any \(\Pi\)-cell \(\sigma\). Use
its rational vertex/recession representation (8), \[
 x=\sum_i\beta_iv_i+\sum_jt_jr_j,\qquad
 \beta_i,t_j\geq0,\quad\sum_i\beta_i=1.
\] Select \(i_0\) with \(\beta_{i_0}>0\), set \(T=\sum_jt_j\), and take
integers \(K\to\infty\) with \(T/K\leq\beta_{i_0}\). Replace \(T/K\) of
that vertex coefficient by coefficients \(t_j/K\) on the rational
endpoints \(v_{i_0}+Kr_j\in\sigma\). This is a same-cell PCC with centre
\textbf{exactly} \(x\), with permitted real coefficients. Affinity of
\(P\) on \(\sigma\), the vertex bounds and the upper growth bound at the
far endpoints give \[
             q(x)\leq P(x)+(T/K)\max\{C_q,0\}.               \tag{38}
\] For \(T=0\) the ordinary vertex PCC already suffices. In general let
\(K\to\infty\) to get \(q(x)\leq P(x)\). This treats every bounded edge,
cycle, unbounded end and other component, with actual allowed cofaces
and all real centres. Conversely \(P\) is an actual PC global competitor
below \(P+b_0\), and \(P\leq\gamma+O(1)\). Therefore \[
                             P_\gamma(u)=P\quad\hbox{on all }X.       \tag{39}
\] This is the raw envelope equality, not an envelope for a replacement
reference.

\subsection{Complete primitive-normal junction
computation}\label{complete-primitive-normal-junction-computation}

In the original product quadrant \((s v_i,t w_j)\), \(s,t\geq0\), put
\(q_i=1/S_k\) if \(i\in I_k\), \(p_j=1/T_l\) if \(j\in J_l\). Directly
from the three specified tripod directions, \[
 H=c_kq_i s+c_lp_j t\quad(k\ne l),\qquad
 H=c_k|q_i s-p_jt|\quad(k=l).                              \tag{40}
\] Let \(a_i=(v_i,0)\), \(b_j=(0,w_j)\). The actual saturated plane
lattice is \(\mathbb{Z}a_i\oplus\mathbb{Z}b_j\): a vector in this plane
belongs to the direct product ambient lattice exactly when its two
primitive factor-line coefficients are integral. Factor star spans may
be nonsaturated in their ray generators; this does not alter these
actual plane lattices.

Off-group quadrants are affine and have index-one axis cones. For a
same-group pair take a positive integer denominator \(D\), \(Q=Dq_i\),
\(P=Dp_j\), \(g=\gcd(Q,P)\). Its zero kink ray has actual primitive
vector \[
                            d_{ij}=(P/g)a_i+(Q/g)b_j.        \tag{41}
\] The integer coefficients are coprime; primitiveness follows also in
the whole ambient lattice because both factor vectors are primitive. The
two actual saturated cone indices are \[
                    \chi_{a_i,d_{ij}}=Q/g,\qquad
                    \chi_{b_j,d_{ij}}=P/g.                 \tag{42}
\] All top cone weights are \(\Omega_{ij}=ab\lambda_i\mu_j\), inherited
from the original graph product. At a kink, (6) and \(H(d_{ij})=0\) give
the full first weight \[
 D_H(d_{ij})=\Omega_{ij}\left(\frac gQ c_kq_i+\frac gP c_kp_j\right)
                             =2c_k\Omega_{ij}g/D.           \tag{43}
\] No denominator or index has been suppressed.

At an axis \(a_i\), the off-group neighbours are \(b_j\). For same-group
neighbours the representative quotient normal is
\(d_{ij}/(Q/g)=b_j+(p_j/q_i)a_i\). Therefore top balancing gives \[
 \sum\frac{\Omega_{ij}}{\chi}r_{ij}
       =ab\lambda_i\left(\sum_j\mu_jb_j+
                    \sum_{j\in J_k}\mu_j(p_j/q_i)a_i\right)
       =ab\lambda_i S_k a_i.                               \tag{44}
\] Since \(H(a_i)=c_k/S_k\), the off-group first term is
\(ab\lambda_i(4-c_k)\), and the tangential correction in (5)--(6), with
(44), subtracts \(ab\lambda_i c_k\). Thus \[
 D_H(a_i)=ab\lambda_i(4-2c_k),\qquad
 D_H(b_j)=ab\mu_j(4-2c_l).                                 \tag{45}
\] These axis rays and the kink rays (41) are all rays in the local
function fan and all are nonnegative. Its balancing can be checked
directly: kink contributions in the first projection for a fixed
\(i\in I_k\) sum to \(2c_kab\lambda_i v_i\), adding its axis
contribution gives \(4ab\lambda_i v_i\). The second projection similarly
gives \(4ab\mu_jw_j\). Both total sums vanish by the original relations.
This check retains every same-group pair even when multiple original
rays share a projected direction.

Apply the second divisor (7). Kink rays have zero value. The entire
first-axis contribution is \[
 ab\sum_k(4-2c_k)c_k\sum_{i\in I_k}\lambda_i/S_k=4ab,
\] and the second-axis contribution is \(4ab\). Hence the
\textbf{complete} coefficient at the chosen junction is \[
                          (H^2\cdot[X])\{(A,B)\}=8ab.              \tag{46}
\] This is the two-stage computation at the original singular product
junction, not a pushed-forward tripod calculation or a one apparent atom
claim.

By (35) the local germ of \(P-P(o)\) is \(\varepsilon H\). Constants
have zero divisor, and both stages of (5) are real linear in the
function, so the complete original-junction coefficient for \(P\) is \[
                    \mathrm{MA}_{\mathrm{poly},\mathbb{R}}(P)\{o\}
                                  =8\varepsilon^2ab.          \tag{47}
\] Other homogeneous cuts introduced by \(\gamma\) or the
carrier/refinement do not change this coefficient, by the proved
two-stage refinement compatibility. The calculation uses no pushforward
to a tripod, no missing-axis convention and no singular-face Alexandrov
substitution.

\subsection{Full strictly negative contact and the fixed-real-pair
conclusion}\label{full-strictly-negative-contact-and-the-fixed-real-pair-conclusion}

The complete PA measure is a finite nonnegative zero-cycle on
\(\Pi^{(0)}\). The exact gap from (37) and (39) is \(P-\gamma-u=-b_0\).
It is \(-1\) at \(o\) and zero at every other \(\Pi\)-vertex. The full
integral therefore is \[
 \int_X(P_\gamma(u)-\gamma-u)
            \mathrm{MA}_{\mathrm{poly},\mathbb{R}}(P_\gamma(u))
                       =-8\varepsilon^2ab<0.                 \tag{48}
\] Refining for the obstacle may introduce vertices with nonzero bump.
The proved successive refinement compatibility gives them zero second
coefficient, since they were not \(\Pi\)-vertices. All original lower
faces, axes, kinks, bounded cycles, ends and disconnected components are
included. Definition 4.4's PA operator is this finite zero-cycle on
\(X\); it adds no point at infinity. The decreasing uniform
regularization \(P+1/n\) has \textbf{identical} measures, so contributes
no escaping mass or infinity term. Equality of total masses was not used
to infer contact.

The decomposition is exact for \textbf{every fixed real reference},
every component \(\gamma_j\) has strict growth on every recession
generator, and the finitely many localization sets are compact. Hence
the reference remains unchanged at infinity, including all other
component pairs. A bad component pair suffices when both factors have a
branch: choose its junctions for the construction; the global norm,
growth estimate, raw vertex/recession proof and zero bump on the other
components apply to the entire original product. If one factor is
unbranched, the reconstructed cylinder proof gives (2) for every bounded
continuous real-PA obstacle and every prescribed real reference. This
proves the iff exactly in the stated finite two-graph subclass,
retaining all the specified boundary and arithmetic cases.

The original arbitrary-balanced-space P05.2 fixed-pair classification
remains open. The fixed-complex cone lemma is not an envelope regularity
theorem; the contact mechanism and the negative mass calculation use the
specific two-graph geometry. No claim is made about the catalogue's
remaining combined intrinsic smoothness/strict-positivity package or the
other residual scopes.

\subsection{Review and reproducibility}

The complete frozen argument reconstructs the finite rational cone of
actual polyhedral convex functions, the exact positive-growth convex
decomposition, the graph-cylinder contact proof, the global raw envelope,
and the full two-stage intersection calculation. The independent Xean
native Codex review used \texttt{gpt-6.1-sol} with \texttt{xhigh}
reasoning. Its complete report and exact argument, manuscript and status
bindings were read and checked before publication.

The root read the entire authoritative argument, both canonical native
notes, author audit and research limits. The exact quadratic-field
diagnostic was recomputed in a fresh directory, preserving the original
artifacts, and matched the retained output. It checks one convex
decomposition and its common rational sample-ray walls. The earlier
rational-reference diagnostics retain the actual saturated cone indices.
These computations do not establish the universal envelope or integral
theorem; those are proved above.

The complete frozen argument, authentic source definitions, standalone
source, PDF, independent review and reproducibility packet accompany the
paper. The standalone source was opened in the native LaTeX editor for
compilation. Same-model review, file bindings, computation and typesetting
are not formal verification or priority certification. The original
classification on arbitrary balanced spaces remains open.

Bibliography: A. M. Botero, E. Mazzon and L. Pille-Schneider,
\emph{Monge--Ampère measures on balanced polyhedral spaces},
arXiv:2603.08664v1, manuscript dated 10 March 2026. Definitions
2.1--2.4, 2.18--2.20, 2.25--2.31, Remark 2.37, Definitions 3.1--3.3,
3.14, 3.20, 3.23, 4.4, 4.15 and 6.2 specify the conventions used here.

The exact defining source quotations remain in full in the accompanying
frozen argument. Every load-bearing supporting mathematical proof above
is retained in this typeset copy.

\end{document}
